\chapter{Introduction}
\label{ch:intro}


%
% Section: Motivation
%
\section{Motivation}

The primary motivation for this internship was to deepen technical knowledge in both hardware and software domains through practical, hands-on projects involving quantum technology and embedded systems. Bell’s Inequality, a cornerstone in quantum mechanics, provided an ideal platform for exploring FPGA-based acceleration, as it requires both precision in design and speed in processing. This hardware-centric task expanded understanding of VHDL implementation, finite state machines, and \acs{FPGA} functionality, all of which are critical for high-performance computing tasks within quantum technology. The hands-on experience gained here also enhanced skills in hardware debugging and optimized design strategies, aligning with the industry's demand for efficient quantum experiment systems.

The second phase of the internship focused on developing a user-centric Python \acs{GUI} based on the qudi framework, motivated by the need for accessible and interactive quantum systems. The \acs{GUI} aimed to streamline user control over \acs{quEDU} and \acs{quNV} devices, providing a customizable platform that allows users to implement their own Python programs and easily interact with hardware. These tools are essential for conducting experiments involving nitrogen-vacancy centers in diamond.









